% p10-05 例 2：矩形 AB=4 AD=8。沿 EF 折，E 在 AD，F 在 BC，使 A 与 C 重合
% A(0,0), B(4,0), C(4,8), D(0,8)，E=(0,5), F=(4,3)
\documentclass[tikz,border=4pt]{standalone}
\usepackage{ctex}
\usepackage{amsmath}
\usetikzlibrary{calc}
\begin{document}
\begin{tikzpicture}[scale=0.5, thick]
  \coordinate (A) at (0,0);
  \coordinate (B) at (4,0);
  \coordinate (C) at (4,8);
  \coordinate (D) at (0,8);
  \coordinate (E) at (0,5);
  \coordinate (F) at (4,3);
  \coordinate (O) at (2,4);  % AC 中点

  \draw (A)--(B)--(C)--(D)--cycle;
  \draw[red, thick] (E)--(F);            % 折痕
  \draw[gray, dashed] (A)--(C);          % AC

  \filldraw (E) circle (2pt);
  \filldraw (F) circle (2pt);
  \filldraw (O) circle (2pt);

  % EF ⊥ AC 在 O 处的小方块
  % AC 方向 (1,2)/√5；EF 方向 (4,-2)/√20=(2,-1)/√5；垂直 ✓
  \draw[gray] ($(O)+(0.15,-0.07)$) -- ($(O)+(0.04,-0.32)$) -- ($(O)+(-0.21,-0.20)$);

  \node[below left]  at (A) {$A$};
  \node[below right] at (B) {$B$};
  \node[above right] at (C) {$C$};
  \node[above left]  at (D) {$D$};
  \node[left]  at (E) {$E$};
  \node[right] at (F) {$F$};
  \node[above right, font=\footnotesize] at (O) {$O$};

  \node[below] at ($(A)!0.5!(B)$) {\footnotesize $4$};
  \node[right] at ($(B)!0.5!(C)$) {\footnotesize $8$};

  \node[red, font=\footnotesize, right] at (F) {\;折痕 $EF$};
\end{tikzpicture}
\end{document}
